
\begin{exercise}[subtitle={The normal distribution \Coffeecup\Coffeecup\Coffeecup}]
    \label{ex: density of normal distribution}
    Let \(X \sim \normal{\mu, \sigma^2}\).
    \begin{enumerate}[label=(\alph*)]
        \item \label{it: (normal distribution) density}
        Show that if \(\sigma^2 > 0\), then \(X\) has density
        \[
            f_X(x) = \frac1{\sqrt{2\pi \sigma^2}}
            \exp\bigl(-\tfrac{(x-\mu)^2}{2\sigma^2}\bigr).
        \]

        \item\label{it: (normal distribution) moments}
        Prove that the centered moments are given by
        \[
            \E{(X-\mu)^n}
            = \begin{cases}
                \sigma^n (n-1)!! & n \text{ even}
                \\
                0 & n \text{ odd}
            \end{cases}
        \]
        \textbf{Hint:} Use the series expansion of \(\exp\) and compare coefficients
        with the Taylor expansion of the moment generating function.

        \item\label{it: (normal distribution) mean and variance}
        Show that \(\E{X} = \mu\) and \(\var{X} = \sigma^2\).
    \end{enumerate}
    Let \(X \sim \normal{\mu, \Sigma}\) be a multivariate normal random variable.
    \begin{enumerate}[label=(\alph*), resume]
        \item\label{it: (multivariate normal distribution) density}
        Show that if \(\Sigma\) is strictly positive definite, then \(X\) has density
        \[
            f_X(x)
            = \frac1{\sqrt{(2\pi)^\dims \det(\Sigma)}}
            \exp\Bigl(-\tfrac12 \scp[\big]{x - \mu, \Sigma^{-1}(x - \mu)}\Bigr).
        \]
        \textbf{Hint:} What is the density of \(Y\) in the representation \(X
        \overset{d}= \mu + A Y\) from Definition \ref{def: multivariate normal
        distribution} \ref{it: multivariate standard normal}? It may
        be helpful to prove \(\abs{\det(A^{-1})} = \smash{\frac1{\sqrt{\det(\Sigma)}}}\)
        using the multiplicativity of the determinant.
        
        \item\label{it: (multivariate normal distribution) mean and covariance}
        Prove that \(\E{X} = \mu\) and \(\Cov{X} = \Sigma\).

        \item\label{it: (multivariate normal distribution) non-positive semidefinite}
        Assume that \(\Sigma\) is not positive semidefinite, that is
        \(v \Sigma v < 0\) for some vector \(v\) prove that there is no random
        variable \(X\) with characteristic function
        \[
            \charFct_X(t)
            = \exp\bigl(i \scp{\mu, t} - \tfrac12 \scp{t, \Sigma t}\bigr).
        \]
        \textbf{Hint:} Exercise \ref{ex: properties moment generating functions} \ref{it: (charfct) bounded}.
    \end{enumerate}
\end{exercise}
\begin{solution}
\begin{enumerate}[wide=0pt, noitemsep]
    \item[\ref{it: (normal distribution) density}]
    We show that the distribution of \(Z = \mu + \sigma Y\) is equal to a random variable \(X\)
    with density \(f_X\). Let \(g\) be a bounded continuous function, then we
    have
    \[\begin{aligned}
        \E{g(Z)}
        &= \E{g(\mu + \sigma Y)}
        \\
        &= \int_\real g(\mu + \sigma y) \frac1{\sqrt{2\pi}} \exp\bigl(-\tfrac{y^2}2\bigr) dy
        \\
        \overset{x=\mu + \sigma y}&=
        \int_\real g(x) \frac1{\sqrt{2\pi \sigma^2}} \exp\bigl(-\tfrac{(x-\mu)^2}{2\sigma^2}\bigr) dx
        \qquad\quad ``\tfrac{1}{\sigma} dx =  dy''
        \\
        &= \E{g(X)}.
    \end{aligned}
    \]
    Since this is true for all bounded continuous functions \(g\) we have \(X \overset{d}= Z\).
    Consequently the density of \(Z\) is given by \(f_X\).

    \item[\ref{it: (normal distribution) moments}]
    Since \(X-\mu \sim \normal{0, \sigma^2}\) we may assume without
    loss of generality that \(\mu=0\).
    Using the series representation of \(\exp\) in the
    the moment generating function of the normal distribution we have
    \[
        \MGF_X(t)
        = e^{\frac12 \sigma^2 t^2}
        = \sum_{k=0}^\infty \frac{(\frac12 \sigma^2 t^2)^k}{k!}
        = \sum_{k=0}^\infty \frac{\sigma^{2k} t^{2k}}{2^k k!}
        = \sum_{\substack{n=0 \\ n \text{ even}}}^\infty \frac{\sigma^{n} t^{n}}{2^{\frac{n}{2}} (\frac{n}{2})!}.
    \]
    By comparing the coefficients with the Taylor expansion of the moment
    generating function
    \[
        \MGF_X(t) = \sum_{n=0}^\infty \frac{\MGF_X^{(n)}(0)}{n!} t^n
        = \sum_{n=0}^\infty \frac{\E{X^n}}{n!} t^n,
    \]
    we get for odd \(n\), \(\E{X^n} = \MGF_X^{(n)}(0) = 0\) and for even \(n\)
    \[
        \E{X^n} = \MGF_X^{(n)}(0) = \frac{n!}{2^{\frac{n}{2}} (\frac{n}{2})!} \sigma^n = (n-1)!! \sigma^n.
    \]

    \item[\ref{it: (normal distribution) mean and variance}]
    By \ref{it: (normal distribution) moments} we have
    \[
        \E{X} = \E{X - \mu} + \mu = \mu,
        \quad\text{and}\quad
        \var{X} = \E{(X-\mu)^2} = \sigma^2.
    \]

    \item[\ref{it: (multivariate normal distribution) density}]
    The density of \(Y\) is the product of the densities of the independent
    standard normal components \((Y_1, \ldots, Y_\dims)\) and therefore given by
    \begin{equation}
        \label{eq: density of Y in multivariate normal distribution}
        \density_Y(y) = \prod_{i=1}^\dims \frac1{\sqrt{2\pi}} \exp\bigl(-\tfrac{y_i^2}2\bigr)
        = \frac1{\sqrt{(2\pi)^\dims}} \exp\bigl(-\tfrac{\norm{y}^2}2\bigr).
    \end{equation}
    To calculate the density of \(X = A Y + \mu\) we use a change of variables with
    the linear map \(\phi(y) = A y + \mu\).  Since \(\Sigma = A A^\transpose\)
    is invertible, we have that \(A\) is invertible and
    by multiplicativity of the determinant
    \[
        1 = \det(\identity) = \det(A A^{-1}) = \det(A) \det(A^{-1}).
    \]
    As this implies \(\det(A^{-1}) = \frac1{\det(A)}\) we thus get
    by multiplicativity of the determinant
    \begin{align*}
        \det(\Sigma) = \det(A A^\transpose) = \det(A)^2
        &\implies \abs{\det(A)} = \sqrt{\det(\Sigma)}
        \\
        &\implies \abs{\det(A^{-1})} = \frac1{\sqrt{\det(\Sigma)}}.
    \end{align*}
    The absolute determinant of the Jacobian of \(\phi^{-1}(x) = A^{-1}(x-\mu)\) is therefore given by
    \begin{equation}
        \label{eq: determinant of Jacobian of phi^-1}    
        \abs{\det(D\phi^{-1}(x))} = \abs{\det(A^{-1})} = \frac1{\sqrt{\det(\Sigma)}}.
    \end{equation}
    We therefore get for any bounded continuous functions \(f\)
    \begin{align*}
        \E{f(X)}
        &= \E{f(A Y + \mu)}
        \\
        \overset{\phi(y) = A y + \mu}&=
        \int_{\real^\dims} f\bigl(\phi(y)\bigr) \density_Y\bigl(y\bigr) dy
        \\
        \overset{\phi^{-1}\text{-subst.}}&= \int_{\real^\dims} f(x) \density_Y\bigl(\phi^{-1}(x)\bigr) \abs[\big]{\det(D\phi^{-1}(x))} dx
        \\
        \overset{\eqref{eq: density of Y in multivariate normal distribution}, \eqref{eq: determinant of Jacobian of phi^-1}}&=
        \int_{\real^\dims} f(x) \underbrace{\frac1{\sqrt{(2\pi)^\dims \det(\Sigma)}} \exp\Bigl(-\tfrac12 \norm[\big]{A^{-1}(x-\mu)}\Bigr)}_{=\density_X(x)} dx
    \end{align*}
    we thereby have identified the density \(\density_X\) of \(X\) as the formula above
    holds for all bounded continuous functions \(f\). This density has the desired form since
    \[
        \norm{A^{-1}(x-\mu)}^2
        = (x-\mu)^\transpose \underbrace{(A^{-1})^\transpose A^{-1}}_{=(AA^\transpose)^{-1} \mathrlap{= \Sigma^{-1}}} (x-\mu)
        = \scp[\big]{x-\mu, \Sigma^{-1}(x-\mu)}.
    \]

    \item[\ref{it: (multivariate normal distribution) mean and covariance}]
    By definition we have \(X \overset{d}= \mu + A Y\) and therefore
    \[
        \E{X} = \mu + A \E{Y} = \mu.
    \]
    Since \(X-\E{X} \overset{d} AY\) we moreover have
    \begin{align*}
        \Cov{X} = \E{(X-\E{X})(X-\E{X})^\transpose}
        &= \E{(A Y)(A Y)^\transpose}
        \\
        &= A \E{Y Y^\transpose} A^\transpose
        = A A^\transpose
        = \Sigma,
    \end{align*}
    since \(\E{Y_i Y_j} = 0\) for \(i \neq j\) and \(\E{Y_i^2} = 1\)
    such that \(\E{Y Y^\transpose} = \identity\).
    
    \item[\ref{it: (multivariate normal distribution) non-positive semidefinite}]
    Assume that \(\Sigma\) is not positive semidefinite and that there
    does exist a random vector \(X\) with the desired characteristic function.
    Since \(\Sigma\) is not positive semidefinite there exists    a vector \(v\)
    such that \(v^\transpose \Sigma v < 0\). Choosing \(t = \lambda v\) for some
    \(\lambda \gg 1\) we get
    \begin{align*}
        \abs{\charFct_X(t)}
        &= \underbrace{\abs[\big]{\exp\bigl(i\scp{\mu, t}\bigr)}}_{=1}\abs[\big]{\exp\bigl(-\tfrac12 \scp{t, \Sigma t}\bigr)}
        \\
        &= \exp\bigl(-\tfrac12 \scp{t, \Sigma t}\bigr)
        = \exp\bigl(-\tfrac{\lambda^2}2 v^\transpose \Sigma v\bigr)
        \overset{\lambda \to \infty}\to \infty,
    \end{align*}
    which contradicts Exercise \ref{ex: properties moment generating functions} \ref{it: (charfct) bounded}.
    Consequently there cannot exist such a random vector \(X\)
    with this characteristic function.
    \qedhere
\end{enumerate}
\end{solution}
